\begin{abcliste}
\abc
Vertikal starten alle vier Kugeln ohne Geschwindigkeit, fallen gleich hoch und haben dieselbe Beschleunigung $g$: Schwerere Kugeln fallen nicht schneller, und die horizontale Geschwindigkeit ändert die vertikale Bewegung nicht. Also
\[
 \result{t_\mathrm{A} = t_\mathrm{B} = t_\mathrm{C} = t_\mathrm{D}}
\]
Alle haben denselben Rang.

\abc
Horizontal wirkt keine Kraft (Trägheitsprinzip): Jede Kugel behält ihre Geschwindigkeit $v$ und landet $x = v\,t$ vom Tisch entfernt. Da alle gleich lang fliegen, entscheidet allein $v$; die Masse (und auch Masse mal Geschwindigkeit) spielt keine Rolle:
\[
 \result{x_\mathrm{B} > x_\mathrm{C} = x_\mathrm{D} > x_\mathrm{A}}
\]

\abc
Aus $h = \frac12\,g\,t^2$:
\begin{align}
 t &= \tFF = \sqrt{\frac{2\cdot\hT}{\g}} = \tF \approx \result{\tFP{0}{2}}
\end{align}
Die Landeweiten:
\begin{align}
 x_\mathrm{A} &= \xaF = \va\cdot\tF = \xa \approx \result{\xaP{0}{2}} \\
 x_\mathrm{B} &= \xbF = \vb\cdot\tF = \xb \approx \result{\xbP{0}{2}} \\
 x_\mathrm{C} &= \xcF = \vc\cdot\tF = \xc \approx \result{\xcP{0}{2}} \\
 x_\mathrm{D} &= \xdF = \vd\cdot\tF = \xd \approx \result{\xdP{0}{2}}
\end{align}
\end{abcliste}
